\section{Limitations and scope}
\label{sec:limitations}

The partition recursions are exact for the stated offline, contiguous, locally factorized segmentation model when each segment marginal likelihood is exact or bounded. The following conditions define the scope of the method and the empirical conclusions.

\paragraph{Computational regime.}
The exact dynamic program has worst-case cost $\Theta(k_{\max}n^2)$ and is designed for offline analysis of complete sequences. For sequences where $n$ or $k_{\max}$ prohibit full computation, a restricted candidate set, a pruning rule with its own correctness analysis, or an approximate calculation can be used; the finite-range timing results in Section~\ref{sec:results} characterize this implementation over the measured range and do not alter the asymptotic bound.

\paragraph{Candidate changepoints and dependence across segments.}
Only changepoints in the candidate set are recoverable. Nonlocal segment interactions and priors that cannot be expressed by local transition factors fall outside the dynamic program derived here. Some forms of dependence can be incorporated through an enlarged state space, but this requires a separate recursion and analysis.

\paragraph{Observation-model scope.}
Posterior conclusions are conditional on the chosen segment model, descriptor semantics, hyperparameters, and within-segment independence. The method does not itself select the appropriate observation family; that choice should be made on scientific grounds before fitting. The array-CGH analysis uses per-probe precision weights derived from rolling median absolute deviation, treated as known inverse-variance scale factors rather than likelihood power weights; comparison to platform-calibrated or uniform-precision specifications remains a natural extension.

\paragraph{Prior sensitivity.}
Posterior distributions over segment count and changepoint location can be sensitive to $p(k)$, $c_u(i,j)$, and $h_u(j)$, particularly with irregularly spaced observations. Sensitivity analyses over prior parameters should be reported alongside primary results.

\paragraph{Nonconjugate approximation.}
Theorem~\ref{thm:stability-paper} provides posterior-odds perturbation bounds conditional on a uniform bound $\varepsilon$ on the block log-score error over all reachable segments. A numerical method that establishes this bound yields an error-controlled approximation; one that does not yields an exact dynamic-programming calculation for its approximate score matrix. Among the methods evaluated in Section~\ref{sec:results}, the Laplace approximation achieved a maximum absolute block-score error of $0.038$ across all tested cells, establishing practical control in those settings. Expectation propagation did not satisfy the bound and should not be used for applications requiring certified posterior accuracy.

\paragraph{Latent-group non-uniqueness.}
The finite latent-group criterion is invariant to label permutation and may admit duplicate or collapsed templates under weak separation. Multiple restarts mitigate but do not eliminate local optima. All reported trajectories were monotone, and all restarts were valid; results should be interpreted as finite-sample evidence for the algorithm's behavior in the declared simulation cells, not as identifiability guarantees for the generative mixture.

\paragraph{Shared-changepoint assumptions.}
Exact pooling requires aligned candidate coordinates and conditional sequence independence given the common partition. When sequences have genuinely sequence-specific changepoints or are otherwise dependent, a common-partition model may be inappropriate; the experimental results show that pooling correctly degrades rather than silently producing invalid boundaries under declared misspecification.

\paragraph{Posterior prediction.}
The observation family and coordinate support are constitutive parts of the predictive model. Prediction outside the fitted coordinate range requires an explicit extrapolation model. The predictive score for held-out observations must use the same observation family as the fitted model; substituting a different distribution produces a finite but invalid score.

\paragraph{Evaluation of changepoint accuracy.}
Boundary $\mathrm{F_1}$ depends on the tolerance and matching rule. Agreement with the joint-MAP boundary set is not accuracy against independently annotated truth. The real-data case studies do not have externally verified boundary annotations; the reported boundaries are posterior summaries of the fitted model, not causal attributions to geological, genomic, or financial events.

\BayesBreak is a generalized hierarchical Bayesian method for offline segmentation of ordered data under locally factorized segment models; it is not a universal changepoint procedure. Online detection, changes in slope or curvature, reversible-jump posterior simulation, and segmentation of unordered or graph-structured data require different models or additional theory \citep{adamsmackay2007bocpd,ruggieri2014bcpvs,green1995rjMCMC,fearnhead2006exact}.
